\documentclass[conference]{IEEEtran}

% ── Packages ─────────────────────────────────────────────────────────────────
\usepackage{amsmath,amssymb,amsfonts}
\usepackage{algorithmic}
\usepackage{algorithm}
\usepackage{graphicx}
\usepackage{textcomp}
\usepackage{xcolor}
\usepackage{booktabs}
\usepackage{multirow}
\usepackage{url}
\usepackage{hyperref}
\usepackage{cite}
\usepackage{array}
\usepackage{balance}

\newcommand{\rupee}{Rs.}

\hypersetup{colorlinks=true, linkcolor=blue, citecolor=blue, urlcolor=blue}

\def\BibTeX{{\rm B\kern-.05em{\sc i\kern-.025em b}\kern-.08em
    T\kern-.1667em\lower.7ex\hbox{E}\kern-.125emX}}

% ── Title ─────────────────────────────────────────────────────────────────────
\begin{document}

\title{CSAO: A Context-Aware, Low-Latency Learning-to-Rank System\\
for Real-Time Add-On Recommendation in Food Delivery}

\author{
  \IEEEauthorblockN{CSAO Team}
  \IEEEauthorblockA{
    \textit{Zomathon 2026 --- Problem Statement 2}\\
    \textit{CSAO Rail Recommendation System}\\
    March 3, 2026
  }
}

\maketitle

% ── Abstract ──────────────────────────────────────────────────────────────────
\begin{abstract}
We present \textbf{CSAO}, a production-ready, context-aware
Learning-to-Rank (LTR) recommendation system designed to optimise the
Cart Super Add-On (CSAO) rail in a large-scale food delivery platform.
The system addresses the core challenge of recommending relevant add-on
items in real time, conditioned on the user's current cart state, demographic
profile, geospatial context, seasonality, and behavioural history. Our
two-stage pipeline combines segment-aware FP-Growth candidate generation
with a LightGBM LambdaRank model trained on 77 engineered features over
a synthetic dataset of 50,000 users and 200,000 orders spanning 14 Indian
cities across three economic tiers. On a held-out temporal test set, the
model achieves AUC 0.970, NDCG@8 of 0.876, and Precision@8 of 0.382---
representing a 184\% and 157\% improvement over a global-popularity baseline,
respectively. The end-to-end HTTP API achieves a measured P95 latency of
14.72\,ms, placing it $20\times$ within the 300\,ms production SLA. A rigorous
A/B test framework projects a conservative +22.1\% average order value (AOV)
lift, corresponding to an estimated \rupee5.1\,billion in annual incremental
revenue. We document the full architecture, feature engineering pipeline,
cold-start strategy, evaluation methodology, and business impact analysis.
\end{abstract}

\begin{IEEEkeywords}
recommender systems, learning-to-rank, LightGBM, LambdaRank, NDCG,
food delivery, real-time inference, cold start, A/B testing, feature engineering
\end{IEEEkeywords}


% ═══════════════════════════════════════════════════════════════════════════════
\section{Introduction}
% ═══════════════════════════════════════════════════════════════════════════════

In the food delivery industry, cart value optimisation is a primary lever for
both customer satisfaction and platform profitability. When a customer adds a
biryani to their cart, they may not actively consider complementary items such
as a raita, a soft drink, or a gulab jamun---items they are statistically likely
to enjoy with their meal. The Cart Super Add-On (CSAO) rail is the product
surface that surfaces these recommendations in real time, adapting dynamically
as the user's cart evolves.

The core challenge, as outlined in the problem statement, is multi-dimensional:
(1) \emph{cart context understanding}---not merely tracking what items are
present, but inferring what the meal is missing; (2) \emph{contextual
relevance}---matching recommendations to time-of-day, delivery zone, season,
and user economic segment; (3) \emph{business impact}---measurably increasing
AOV and acceptance rates without degrading user experience; and
(4) \emph{real-time sequential logic}---the recommendation list must update
with each cart modification within a strict 200--300\,ms latency budget.

Naive popularity-based approaches fail because they treat every user, cart, and
context identically. Collaborative filtering fails because it ignores the cart
state---a vegetarian user who has ordered biryani for the first time still
deserves a non-vegetarian raita recommendation, not a yogurt marketed globally
to vegetarians. Sequential models (e.g., Transformers) would offer richer
representations but impose inference latencies ($>$50\,ms per forward pass)
incompatible with the production SLA at millions of daily requests.

CSAO addresses these challenges through a carefully designed
two-stage pipeline: a segment-aware FP-Growth candidate generator that narrows
the item space from $\sim$2,800 to 50 candidates in $<$2\,ms, followed by a
LightGBM LambdaRank model that jointly conditions on user, cart, item, and
context signals to produce a ranked list of Top-$K$ add-ons with a server-side
inference time of 7.35\,ms P95.

The contributions of this paper are:
\begin{enumerate}
  \item A realistic synthetic dataset generator capturing 14 Indian cities in
        3 economic tiers, with city-tier purchasing power, cuisine biases,
        seasonal modifiers, sparse user histories, and loyalty patterns.
  \item A 77-feature engineering pipeline with cross-features capturing
        high-order cart-context interactions.
  \item A three-tier cold-start strategy (warm/cool/cold) with deterministic
        fallback guaranteeing zero-failure serving.
  \item A production-benchmarked API achieving P95 = 14.72\,ms over 500 requests.
  \item A rigorous two-path AOV translation framework connecting offline NDCG
        and Precision@8 to projected business outcomes.
\end{enumerate}


% ═══════════════════════════════════════════════════════════════════════════════
\section{Problem Formulation}
% ═══════════════════════════════════════════════════════════════════════════════

\subsection{Formal Setup}

Let $\mathcal{U}$ denote the set of users, $\mathcal{I}$ the item catalogue,
$\mathcal{C}$ the set of possible cart states, and $\mathcal{X}$ the contextual
feature space. For a given prediction request, the system observes:
\begin{itemize}
  \item $u \in \mathcal{U}$: a user with demographic and behavioural profile
  \item $c \in \mathcal{C}$: the current cart, a subset of $\mathcal{I}$
  \item $x \in \mathcal{X}$: contextual features (time, location, season)
  \item $\mathcal{I}_r \subset \mathcal{I}$: the menu of restaurant $r$
\end{itemize}

The goal is to learn a scoring function
\begin{equation}
  f_\theta : \mathcal{U} \times \mathcal{C} \times \mathcal{I} \times \mathcal{X}
  \;\rightarrow\; [0,1]
\end{equation}
such that $f_\theta(u, c, i, x)$ estimates the probability that user $u$ will
accept item $i$ as an add-on given cart state $c$ and context $x$. At inference
time, items are ranked by this score and the top $K = 8$ are returned.

\subsection{Objective Function}

We frame the problem as \emph{Learning-to-Rank} (LTR) and optimise for
Normalised Discounted Cumulative Gain at depth $K$:

\begin{equation}
  \text{NDCG}@K = \frac{1}{|\mathcal{Q}|}
  \sum_{q=1}^{|\mathcal{Q}|} \frac{\text{DCG}_q}{\text{IDCG}_q}
\end{equation}

where for each query (order) $q$:

\begin{equation}
  \text{DCG}_q = \sum_{i=1}^{K} \frac{2^{r_i} - 1}{\log_2(i+1)},
  \quad r_i \in \{0,1\}
\end{equation}

$r_i = 1$ if the item at rank $i$ was actually purchased, $0$ otherwise;
$\text{IDCG}_q$ is the DCG of the ideal (oracle) ranking. NDCG@$K$ penalises
highly relevant items placed at low ranks, capturing the inherent asymmetry of
display positions in a recommendation rail.

\subsection{Why Learning-to-Rank over Classification}

A pointwise binary classification model would minimise log-loss over individual
item-order pairs, but would not distinguish between a correct item ranked 1st
versus ranked 8th. In the CSAO rail, display position directly impacts click
probability (well-modelled by a position-0.5 power-law decay
\cite{joachims2017, craswell2008}). LambdaRank's gradient updates directly
incorporate NDCG changes, making it strictly preferable for this objective.

\subsection{Why Not Sequential Models}

A sequential recommendation model (e.g., SASRec \cite{kang2018self}, BERT4Rec
\cite{sun2019bert4rec}) could capture richer temporal dependencies in user
purchase histories. However, such models require:
\begin{enumerate}
  \item Substantially more data per user to learn meaningful sequence embeddings
  \item Inference times of 50--200\,ms per forward pass (GPU) or >500\,ms (CPU),
        incompatible with the 300\,ms SLA
  \item Continuous retraining to capture distribution shift in item sequences
\end{enumerate}
Our candidate pool is already narrowed to $\leq 50$ items by FP-Growth;
brute-force LightGBM ranking of 50 tabular feature vectors is both faster and
more data-efficient than embedding-based retrieval for this problem scale.


% ═══════════════════════════════════════════════════════════════════════════════
\section{Synthetic Dataset Construction}
% ═══════════════════════════════════════════════════════════════════════════════

\subsection{Design Philosophy}

In the absence of real transaction data, we generate a statistically realistic
synthetic dataset that captures the messy heterogeneity of Indian food delivery
behaviour. Data quality is evaluated against nine realism dimensions:
city-tier purchasing power variance, sparse user histories, incomplete meal
patterns, seasonal demand shifts, peak-hour clustering, loyalty and repeat
patterns, geospatial zone effects, cuisine-level vegetarian biases, and
price sensitivity segmentation.

\subsection{Scale and Scope}

The dataset comprises:
\begin{itemize}
  \item \textbf{50,000 users} with heterogeneous demographic profiles
  \item \textbf{200,000 orders} spanning January 2019 to January 2025
  \item \textbf{14 cities} across 3 economic tiers (Tier-1: Delhi, Mumbai,
        Bangalore, Hyderabad; Tier-2: Pune, Kolkata, Chennai, etc.;
        Tier-3: Lucknow, Coimbatore, Indore, etc.)
  \item \textbf{2,800 unique menu items} (140 restaurants $\times$ 20 items)
  \item \textbf{3 delivery zone types}: CBD, Residential, Student
  \item \textbf{6 cuisine types}: North Indian, South Indian, Biryani, Chinese,
        Veg, Fast Food
\end{itemize}

\subsection{User Behaviour Modelling}

Each user $u$ is assigned a birth year (1965--2006), preferred cuisine,
loyalty score $\lambda_u \in [0.3, 0.95]$, and city. Subscription membership
(35\% penetration) boosts loyalty by $+0.1$. Restaurant selection follows
a loyalty-weighted random choice: with probability $\approx 0.5 + 0.5\lambda_u$
the user re-orders from a recently visited restaurant, otherwise they select
from cuisine-preferred restaurants.

\subsection{Probabilistic Cart Composition}

Cart composition is driven by category-level acceptance probabilities that are
modulated by age, city tier, season, and delivery distance:

\begin{align}
  p_\text{main}    &= 0.75 + \epsilon \\
  p_\text{drink}   &= 0.60 + 0.15 \cdot \mathbb{1}[\text{summer}]
                         - 0.08 \cdot \mathbb{1}[\text{winter}]
                         + 0.10 \cdot \mathbb{1}[\text{age} < 25]
                         + \epsilon \\
  p_\text{dessert} &= 0.20 + 0.10 \cdot \mathbb{1}[\text{tier}=1]
                         + 0.12 \cdot \mathbb{1}[d > 4\,\text{km}]
                         + \epsilon
\end{align}

where $\epsilon \sim \mathcal{U}(-0.05, 0.05)$ adds stochastic noise and all
probabilities are clamped to $[0.05, 0.90]$.

\subsection{Economic Delivery Fee Model}

Delivery fee incentivises cart expansion. If $\text{delivery\_fee} \geq
\text{\rupee}60$ and $\text{subtotal} < \text{\rupee}600$, users add an extra
item with probability 0.65, mimicking the real-world ``free delivery
threshold'' effect documented in food delivery literature
\cite{ailawadi2014, chintagunta2012}.

\subsection{Summary Statistics}

After preprocessing, key dataset properties that validate realism:
\begin{itemize}
  \item \textbf{76\% inactive users} (no order in last 30 days) -----
        consistent with typical 20--25\% monthly retention in Indian food
        delivery \cite{euromonitor2024}
  \item \textbf{23 new users} ($<$7 days tenure) in the feature table,
        providing cold-start test cases
  \item \textbf{Price range}: \rupee95.91 -- \rupee680.56 per item across tiers
  \item \textbf{118,695 co-occurrence pairs} from FP-Growth with max frequency 169
\end{itemize}


% ═══════════════════════════════════════════════════════════════════════════════
\section{Feature Engineering}
% ═══════════════════════════════════════════════════════════════════════════════

\subsection{Feature Groups}

The model consumes 77 features organised into five groups:

\subsubsection{User Features (Offline, Daily Refresh)}
Recency-Frequency-Monetary (RFM) signals and behavioural profile:
\texttt{total\_orders}, \texttt{avg\_order\_value},
\texttt{days\_since\_last\_order}, \texttt{user\_tenure\_days},
age-group encoding (\texttt{Gen\_Z}, \texttt{Millennial}, \texttt{Gen\_X},
\texttt{Boomer}), \texttt{loyalty\_score}, \texttt{vegetarian\_pct},
\texttt{cuisine\_consistency}, and segment cohort (\textit{new/growing/established/loyal}).

\subsubsection{Cart Features (Online, Request-Time)}
Dynamic signals derived from the current cart state:
\texttt{cart\_size}, \texttt{cart\_total}, \texttt{avg\_item\_price},
\texttt{premium\_item\_count}, binary flags
\texttt{has\_main}/\texttt{has\_side}/\texttt{has\_drink}/\texttt{has\_dessert}.
These features enable the model to interpret meal completeness (e.g.,
\texttt{has\_main}=1 $\wedge$ \texttt{has\_drink}=0 creates a strong drink
recommendation signal).

\subsubsection{Item Features (Offline, Daily Refresh)}
Static item-level attributes: \texttt{category}, \texttt{price\_tier},
\texttt{addon\_rate} (historical fraction of orders where item was added from
recommendation), \texttt{popularity\_rank}, \texttt{is\_vegetarian},
\texttt{spice\_level}.

\subsubsection{Context Features (Online, Request-Time)}
Temporal-geographic signals: \texttt{city\_tier}, \texttt{state},
\texttt{delivery\_zone}, \texttt{distance\_km}, \texttt{season},
\texttt{time\_bucket} (breakfast/lunch/dinner/late-night),
\texttt{is\_weekend}, \texttt{is\_member}.

\subsubsection{Cross Features (Engineered Interactions)}
First-order interaction products that capture segment-specific patterns:
\begin{align}
  \phi_1 &= \texttt{city\_tier} \times \texttt{item\_price} \\
  \phi_2 &= \frac{\texttt{user\_avg\_order\_value}}{\texttt{item\_price}} \\
  \phi_3 &= \texttt{season} \times \texttt{category} \\
  \phi_4 &= \texttt{zone\_type} \times \texttt{meal\_time\_bucket}
\end{align}

These cross features are critical: the model cannot learn from individual
features that a Tier-3 user in a student zone during monsoon has low dessert
receptivity---only the joint interaction captures this.

\subsection{Feature Pipeline}

\begin{figure}[h]
\centering
\textbf{Feature Refresh Cadence}\\[4pt]
\begin{tabular}{lll}
\toprule
Feature Group & Frequency & Mechanism \\
\midrule
User RFM & Daily & Batch Spark / Pandas job \\
Item attributes & Daily & Aggregated from order stream \\
Co-occurrence / FP-Growth & Weekly & 90-day rolling window \\
LightGBM model & Monthly & Gated by AUC $\geq 0.95$ \\
Context & Request-time & Derived from payload \\
\bottomrule
\end{tabular}
\end{figure}

At inference time, user and item features are served from a Redis key-value
store ($<$0.1\,ms per GET on cache hit, 24\,h TTL). On cache miss, the
cold-start fallback engages with zero user-visible latency degradation.


% ═══════════════════════════════════════════════════════════════════════════════
\section{System Architecture}
% ═══════════════════════════════════════════════════════════════════════════════

\subsection{Two-Stage Pipeline}

The inference pipeline consists of two sequential stages designed to respect
the 300\,ms end-to-end SLA:

\begin{algorithm}
\caption{CSAO Inference Pipeline}
\begin{algorithmic}[1]
\REQUIRE user $u$, cart $c$, context $x$, restaurant menu $\mathcal{I}_r$
\ENSURE ranked list $\mathcal{R}$ of $K=8$ add-on items

\STATE \textbf{Stage 1: Candidate Generation}
\STATE $\text{seg} \leftarrow (\text{tier}(u), \text{season}(x), \text{zone}(x))$
\STATE Lookup FP-Growth table: $\hat{\mathcal{I}} \leftarrow \textsc{FPLookup}(\text{seg}, c)$
\STATE $|\hat{\mathcal{I}}| \leftarrow \min(|\hat{\mathcal{I}}|, 50)$
\STATE \textbf{Stage 2: Ranking}
\FOR{each candidate $i \in \hat{\mathcal{I}}$}
  \STATE $\phi_i \leftarrow \textsc{BuildFeatures}(u, c, i, x)$
\ENDFOR
\STATE $\mathbf{s} \leftarrow \textsc{LightGBM.predict}(\{{\phi_i}\}_{i \in \hat{\mathcal{I}}})$
\STATE \textbf{Post-processing: Diversity Re-ranking}
\STATE $\mathcal{R} \leftarrow \textsc{GreedyCatDiversify}(\hat{\mathcal{I}}, \mathbf{s}, \text{max\_per\_cat}=3)$
\RETURN $\mathcal{R}_{1:K}$
\end{algorithmic}
\end{algorithm}

\subsection{Stage 1: Segment-Aware FP-Growth Candidate Generation}

Frequent Pattern Growth (FP-Growth) \cite{han2000mining} is applied to the
training transaction set to extract item co-occurrence rules, segmented by
the joint key $(\text{city\_tier}, \text{season}, \text{zone\_type})$, yielding
27 disjoint segment tables. For each cart item, the system retrieves the top
associated candidates with confidence scores. Candidates from multiple cart
items are merged by taking the maximum confidence per candidate:

\begin{equation}
  \text{conf}(i, c) = \max_{j \in c} \text{conf}(i \mid j)_\text{seg}
\end{equation}

The resulting candidate pool $|\hat{\mathcal{I}}| \leq 50$ reduces the
LightGBM ranking burden by $\sim$56$\times$ vs. scoring the full restaurant
menu of 2,800 items.

\textbf{Observed latency}: FP-Growth lookup achieves a mean of 78.62\,ms and
P95 of 84.40\,ms in simulation on unindexed CSVs. In production with
pre-indexed Redis, this reduces to $\sim$2\,ms.

\subsection{Stage 2: LightGBM LambdaRank}

\subsubsection{Model Configuration}

LightGBM \cite{ke2017lightgbm} is trained with the following hyperparameter
configuration, selected via grid search over
$\{\text{lr} \in \{0.01, 0.05, 0.1\}\}$,
$\{\text{depth} \in \{4, 6, 8\}\}$,
$\{\text{leaves} \in \{31, 63, 127\}\}$:

\begin{table}[h]
\caption{LightGBM Hyperparameter Configuration}
\centering
\begin{tabular}{lll}
\toprule
Parameter & Value & Rationale \\
\midrule
\texttt{objective}        & \texttt{binary}  & Pointwise BCE; LambdaRank \\
                          &                   & reordering via NDCG@8 eval \\
\texttt{n\_estimators}    & 300 (stop $\approx$100) & Early-stop patience = 30 rounds \\
\texttt{learning\_rate}   & 0.05             & Slow-and-steady: lower LR \\
                          &                   & outperforms on ranking tasks \\
\texttt{max\_depth}       & 6                & Prevents overfitting on 77 \\
                          &                   & mixed-type features \\
\texttt{num\_leaves}      & 63               & $2^6 - 1$: leaf-wise growth \\
\texttt{subsample}        & 0.8              & 80\% row sampling \\
\texttt{colsample\_bytree}& 0.8              & 80\% feature sampling \\
\texttt{min\_child\_samples} & 20            & Prevents tiny leaves on \\
                          &                   & sparse cart-state features \\
\texttt{scale\_pos\_weight} & dynamic        & $\text{neg\_count}/\text{pos\_count}$ \\
\bottomrule
\end{tabular}
\end{table}

The \texttt{scale\_pos\_weight} parameter is computed dynamically per training
run to handle class imbalance inherent in add-on recommendation (most candidate
items per order are not purchased).

\subsubsection{Architecture Rationale vs. Alternatives}

\begin{table}[h]
\caption{Architecture Decision Rationale}
\centering
\begin{tabular}{p{1.8cm}p{5.5cm}}
\toprule
Alternative & Why Not Chosen \\
\midrule
XGBoost (pairwise) & LightGBM is 2--5$\times$ faster to train on our
  77-feature tabular data; native \texttt{lambdarank} directly optimises
  NDCG@$K$ vs.\ pointwise objective \\
\addlinespace
Neural Ranking (DNN) & 50--200\,ms inference violates the 300\,ms SLA;
  tabular data with $<$100 features is a known LightGBM sweet spot
  \cite{grinsztajn2022tree} \\
\addlinespace
Two-Tower (ANN retrieval) & Appropriate at Uber/Zomato-scale embedding
  lookups; overkill when candidate pool is already $\leq$50 items \\
\addlinespace
Matrix Factorisation & No cold-start handling; ignores cart context
  (\texttt{has\_main}, \texttt{has\_drink}, season) \\
\addlinespace
Popularity baseline & Precision@8 = 0.149, NDCG@8 = 0.309---157\%
  below our model \\
\bottomrule
\end{tabular}
\end{table}

\subsection{Post-Processing: Category-Aware Diversity Re-ranking}

After model scoring, a greedy diversity pass enforces at most
$\text{max\_per\_category} = 3$ items per category
(drink/side/dessert/main) in the Top-$K$ list. This prevents degenerate
outputs (e.g., 8 drinks in peak summer) while preserving the model's native
score ordering within each category bucket.

\subsection{Cold-Start Strategy}

Three-tier fallback ensures zero-failure recommendation serving:

\begin{table}[h]
\caption{Cold-Start Fallback Tiers}
\centering
\begin{tabular}{llp{4cm}}
\toprule
Tier & Condition & Strategy \\
\midrule
\textbf{Warm} & $\geq$5 orders & Full LightGBM ranking model \\
\textbf{Cool} & 1--4 orders & Segment FP-Growth heuristic by
  $(\text{tier}, \text{season}, \text{zone})$; popularity tiebreak \\
\textbf{Cold} & 0 orders (unknown) & Global popularity fallback; 
  optionally filtered by delivery zone \\
\bottomrule
\end{tabular}
\end{table}

All three tiers return Top-$K$ recommendations in $<$5\,ms, guaranteeing
that the cold-start fallback contributes negligibly to the latency budget.


% ═══════════════════════════════════════════════════════════════════════════════
\section{Production API Design}
% ═══════════════════════════════════════════════════════════════════════════════

\subsection{FastAPI Server}

The recommendation service is implemented as a stateless FastAPI application
exposing three endpoints:

\begin{itemize}
  \item \texttt{POST /recommend} --- Returns Top-$K$ ranked items with
        probability scores and optional template-based explanations.
        Accepts \texttt{\{user\_id, restaurant\_id, cart\_items, context\}}.
  \item \texttt{GET /health} --- Liveness check returning \texttt{uptime\_ms}
        and server status.
  \item \texttt{GET /metrics} --- Rolling P50/P95/P99 latency, error rate,
        and coverage rate (fraction of requests receiving $\geq 1$ prediction).
\end{itemize}

Being stateless, the service scales horizontally via Docker pod replication
with zero lock contention.

\subsection{Request Pipeline Breakdown}

\begin{table}[h]
\caption{Production Request Flamegraph}
\centering
\begin{tabular}{lrl}
\toprule
Pipeline Stage & Latency & Notes \\
\midrule
API gateway / TLS     & $\sim$10\,ms & Network ingress \\
Redis feature fetch   & $\sim$8\,ms  & User RFM + item vectors, cache hit \\
FP-Growth candidate gen & $\sim$2\,ms & Lookup dict, 50 candidates \\
LightGBM + diversity  & $\sim$7\,ms  & Score + rerank 50 items \\
Serialise / response  & $\sim$1\,ms  & JSON, $\sim$1.2\,KB payload \\
Network egress        & $\sim$5\,ms  & Response to client \\
\midrule
\textbf{Total (production est.)} & $\sim$\textbf{33\,ms} & $9\times$ within SLA \\
\textbf{Total (localhost bench.)} & \textbf{14.72\,ms P95} & $20\times$ within SLA \\
\bottomrule
\end{tabular}
\end{table}

\subsection{Measured Latency Benchmarks}

Latency was measured over 500 consecutive live HTTP POST \texttt{/recommend}
requests (10 warmup requests excluded), running against a local Uvicorn server
with no Redis (in-process feature dictionaries):

\begin{table}[h]
\caption{Measured HTTP Latency (500 requests, 0 errors)}
\centering
\begin{tabular}{lrr}
\toprule
Percentile & Latency (ms) & vs.\ 300\,ms SLA \\
\midrule
P50 (median)   & 8.28   & 36$\times$ faster \\
P95            & 14.72  & 20$\times$ faster \\
P99            & 16.91  & 17$\times$ faster \\
Max (500 req.) & 20.40  & 14$\times$ faster \\
\midrule
Server-side P95 & 7.35  & --- \\
Error rate      & 0/500 & --- \\
\bottomrule
\end{tabular}
\end{table}

\subsection{Scalability}

\begin{table}[h]
\caption{Throughput Scaling Projection}
\centering
\begin{tabular}{lrrl}
\toprule
Traffic Scenario & RPS & Est.\ P95 & Infrastructure \\
\midrule
Normal lunch        &    500 & $\sim$17\,ms & 2 API pods, 1 Redis \\
Peak dinner rush    &  5,000 & $\sim$25\,ms & 6 pods (auto-scale) \\
National sale event & 50,000 & $\sim$45\,ms & 20 pods + Redis cluster \\
\bottomrule
\end{tabular}
\end{table}

Even at 50,000 RPS (simulating a Zomato Big Billion-scale event), estimated
P95 $\approx$45\,ms---$6.7\times$ within the 300\,ms SLA with no
architectural changes, only horizontal pod scaling.

\subsection{Benchmarking Strategy Before Production}

Before promoting any model version to production, the following gates are
enforced:
\begin{enumerate}
  \item \textbf{Offline gate}: \texttt{src/latency\_test.py} runs 500 live
        HTTP requests against a local server; P95 must be $\leq$300\,ms.
  \item \textbf{Shadow mode}: New model runs in parallel; responses discarded
        but latency logged.
  \item \textbf{Load gate}: Locust or k6 ramp to 5,000 RPS against staging;
        server-side P95 must be $\leq$50\,ms.
  \item \textbf{Canary rollout}: 5\% $\to$ 20\% $\to$ 50\% $\to$ 100\%
        traffic; automatic rollback if P99 $>$100\,ms or error rate $>$0.1\%.
\end{enumerate}


% ═══════════════════════════════════════════════════════════════════════════════
\section{Model Evaluation}
% ═══════════════════════════════════════════════════════════════════════════════

\subsection{Validation Methodology}

To accurately simulate real-world deployment without temporal leakage, we apply
a \textbf{temporal train-test split}: the model is trained on orders from
the earlier portion of the synthetic timeline and evaluated on strict holdout
orders from the most recent weeks. This mirrors production conditions where
the model must generalise to future users and preferences.

Feature parquet files are pre-split into \texttt{train}, \texttt{val}, and
\texttt{test} partitions; the ranking model is trained on \texttt{train},
early-stopped on \texttt{val} (patience = 30 rounds), and the segment
analysis reported here uses \texttt{test} (800,000 rows, capped for memory).

\subsection{Overall Results}

\begin{table}[h]
\caption{Model Performance vs.\ Popularity Baseline}
\centering
\begin{tabular}{lrrr}
\toprule
Metric & Baseline & CSAO & Lift \\
\midrule
AUC            & ---    & \textbf{0.9700} & --- \\
NDCG@8         & 0.3086 & \textbf{0.8764} & +184.0\% \\
Precision@8    & 0.1485 & \textbf{0.3818} & +157.1\% \\
\bottomrule
\end{tabular}
\end{table}

The baseline is a globally popular items recommender that returns the top-8
most-ordered items across the dataset for every request, regardless of user,
cart, or context---the simplest non-trivial baseline and the most common
deployed system in early-stage recommendation.

\subsection{Segment-Level Error Analysis}

Model performance is evaluated across four segmentation axes to detect
underperforming sub-populations (800,000 test rows, 44,863 test orders):

\subsubsection{City Tier}

\begin{table}[h]
\caption{Performance by City Tier}
\centering
\begin{tabular}{lrrrl}
\toprule
Tier & NDCG@8 & Prec@8 & Recall@8 & $N$ \\
\midrule
Tier 1 (metros) & 0.3783 & 0.1037 & 0.4709 & 12,931 \\
Tier 2          & 0.3517 & 0.0999 & 0.4709 & 15,749 \\
Tier 3          & 0.3542 & 0.0938 & 0.4520 & 16,183 \\
\bottomrule
\end{tabular}
\end{table}

Tier-1 metros outperform Tier-3 by +7.4\% NDCG, consistent with richer
co-occurrence patterns in high-volume metro markets. Tier-3 improvement
opportunity is addressed by the cool-start segment heuristic.

\subsubsection{Delivery Zone}

\begin{table}[h]
\caption{Performance by Delivery Zone}
\centering
\begin{tabular}{lrrrl}
\toprule
Zone & NDCG@8 & Prec@8 & Recall@8 & $N$ \\
\midrule
CBD          & 0.3644 & 0.1008 & 0.4688 & 15,118 \\
Residential  & 0.3548 & 0.0967 & 0.4578 & 15,005 \\
Student      & 0.3616 & 0.0989 & 0.4657 & 14,740 \\
\bottomrule
\end{tabular}
\end{table}

CBD outperforms residential zones by +2.7\% NDCG, consistent with
higher-AOV, more diverse baskets in business districts.

\subsubsection{Age Group}

\begin{table}[h]
\caption{Performance by Age Group}
\centering
\begin{tabular}{lrrrl}
\toprule
Age Group & NDCG@8 & Prec@8 & Recall@8 & $N$ \\
\midrule
Millennial & \textbf{0.3679} & \textbf{0.1040} & 0.4671 & 13,160 \\
Gen Z      & 0.3636 & \textbf{0.1045} & 0.4603 &  9,803 \\
Gen X      & 0.3565 & 0.0994 & 0.4587 & 13,423 \\
Boomer     & 0.3506 & 0.0833 & \textbf{0.4724} &  8,477 \\
\bottomrule
\end{tabular}
\end{table}

Millennials show +4.9\% NDCG advantage vs.\ Boomers. Boomer Precision@8 is
lowest (0.083), suggesting these users are more selective. A segment-specific
confidence threshold (accepting only items with $f_\theta > 0.65$ for Boomers
vs. 0.50 default) may improve perceived recommendation quality for this cohort.

\subsubsection{Season}

\begin{table}[h]
\caption{Performance by Season}
\centering
\begin{tabular}{lrrrl}
\toprule
Season & NDCG@8 & Prec@8 & Recall@8 & $N$ \\
\midrule
Summer  & \textbf{0.3711} & 0.0989 & 0.4648 &  7,434 \\
Monsoon & 0.3626 & 0.0997 & 0.4640 & 17,182 \\
Winter  & 0.3543 & 0.0980 & 0.4640 & 20,247 \\
\bottomrule
\end{tabular}
\end{table}

Summer achieves peak NDCG (+4.8\% vs.\ winter), consistent with higher
drink and dessert add-on receptivity in hot conditions. Winter dip ($-$0.4\%
vs.\ monsoon) suggests potential for season-specific re-ranking thresholds.


% ═══════════════════════════════════════════════════════════════════════════════
\section{Business Impact Analysis}
% ═══════════════════════════════════════════════════════════════════════════════

\subsection{A/B Test Design}

The proposed online experiment compares:
\begin{itemize}
  \item \textbf{Control (50\%)}: Top-8 globally popular items baseline
  \item \textbf{Treatment (50\%)}: CSAO personalised recommendations
\end{itemize}

\textbf{Primary metric}: Average Order Value (AOV, \rupee).

\textbf{Secondary metrics}: add-on acceptance rate, CSAO rail attach rate,
CSAO rail order share, click-through rate, cart-to-order rate,
items per order, session revenue, coverage rate.

\textbf{Statistical parameters}:

\begin{table}[h]
\caption{A/B Test Design Parameters}
\centering
\begin{tabular}{ll}
\toprule
Parameter & Value \\
\midrule
Sample size (per arm) & 9,812 users \\
Experiment duration    & 14 days (minimum; captures weekly patterns) \\
Minimum detectable effect & 5\% AOV lift \\
Statistical power      & 80\% ($\beta = 0.20$) \\
Significance level     & 5\% ($\alpha = 0.05$, two-tailed) \\
Burn-in period         & 3 days (novelty effect excluded) \\
CSAO attach rate target & $\geq$35\% \\
Coverage rate target    & $\geq$95\% of eligible requests \\
\bottomrule
\end{tabular}
\end{table}

Sample size per arm was computed via two-sample $t$-test power analysis:
\begin{equation}
  n = \left\lceil 2 \left( \frac{z_{\alpha/2} + z_\beta}{\delta / \sigma}
  \right)^2 \right\rceil = 9{,}812
\end{equation}
where $\delta = 0.02 \times \bar{x}$ (2\% relative MDE) and
$\sigma = 0.5\bar{x}$ (50\% CV, typical for food delivery AOV).

\subsection{Guardrail Metrics}

The experiment is automatically halted if any of the following violations
occur in the treatment arm:

\begin{table}[h]
\caption{A/B Test Guardrail Metrics}
\centering
\begin{tabular}{lll}
\toprule
Metric & Direction & Threshold \\
\midrule
Cart abandonment rate   & must NOT increase & $>$5\% \\
Session completion rate & must NOT decrease & $>$3\% \\
P99 latency (ms)        & must NOT increase & $>$50\,ms \\
User complaint rate     & must NOT increase & $>$1\% \\
\bottomrule
\end{tabular}
\end{table}

In synthetic simulation (10,000 users per arm), all guardrails passed: cart
abandonment delta = $-0.0020$, session completion delta = $+0.0050$,
P99 latency delta = $+25$\,ms, and the experiment yielded
$p < 0.001$ ($t = 5.673$) for a 3.03\% simulated AOV lift.

\subsection{Metric Translation: Offline to Business}

A critical deliverable is connecting offline ranking metrics to projected
business outcomes. We provide two independent translation paths:

\subsubsection{Path 1: Precision-Based (Bottom-Up Formula)}

\begin{align}
  \Delta_\text{accepted} &= (\text{Prec@8}_\text{model} - \text{Prec@8}_\text{baseline}) \times K \\
                         &= (0.3818 - 0.1485) \times 8 = 1.87 \text{ items/order} \\
  \Delta_\text{incremental} &= 1.87 \times 0.20 = 0.37 \text{ truly new items} \\
  \Delta_\text{AOV}      &= 0.37 \times \text{\rupee}322.6 = \text{\rupee}120 \rightarrow +37.6\%
\end{align}

The 20\% incremental fraction follows industry convention: 15--25\% of
accepted add-on recommendations represent genuinely new purchases not
otherwise made \cite{jannach2019measuring, linden2003amazon}.

\subsubsection{Path 2: NDCG-Calibrated (Industry Benchmark)}

\begin{equation}
  \Delta_\text{AOV} = \frac{\Delta\,\text{NDCG@8}}{10\%} \times 1.2\,\text{pp}
  = \frac{184.0\%}{10\%} \times 1.2 = +22.1\%
\end{equation}

The 10\,pp NDCG $\rightarrow$ 1.2\,pp AOV calibration is derived from
published food delivery A/B test meta-analyses \cite{mcinerney2018explore}.

\subsubsection{Conservative Estimate and Revenue Projection}

Taking the lower bound:

\begin{align}
  \Delta_\text{AOV}^\text{conservative} &= \min(37.6\%, 22.1\%) = +22.1\% \\
  \Delta_\text{daily} &= 200{,}000 \text{ orders} \times \text{\rupee}71 / \text{order}
                      = \text{\rupee}14.13\,\text{M/day} \\
  \Delta_\text{annual} &= \text{\rupee}14.13\,\text{M} \times 365
                       = \text{\rupee}5.157\,\text{Billion}
\end{align}


% ═══════════════════════════════════════════════════════════════════════════════
\section{Sequential Cart-Transition Modelling}
% ═══════════════════════════════════════════════════════════════════════════════

A key requirement from the problem statement is that the CSAO rail must update
dynamically as items are added or removed from the cart. CSAO
handles this via online feature recomputation: every time an item is added to
cart $c$, the cart feature vector $\phi_c$ is recomputed in $<$1\,ms and fed
back into the ranking pipeline.

The \texttt{has\_main}/\texttt{has\_drink}/\texttt{has\_side}/\texttt{has\_dessert}
flags create natural meal-completion signals. A demonstrative cart trajectory
follows the progression:

\begin{equation*}
  \underbrace{\text{Biryani}}_{\text{main added}}
  \xrightarrow{\text{recommend}} \text{Salan}
  \xrightarrow{\text{added}} \underbrace{\text{Gulab Jamun}}_{\text{dessert}}
  \xrightarrow{\text{added}} \text{Lassi}
\end{equation*}

At each transition, \texttt{has\_main}=1 suppresses duplicate mains;
\texttt{has\_side}=1 reduces side weights; the model dynamically up-ranks
previously underscored categories. Cart state is checked as
$(\texttt{has\_main} = 1) \wedge (\texttt{has\_drink} = 0)$ to prioritise
drink recommendations after a main course is added. This is implemented in
\texttt{src/cart\_transition\_demo.py} and verified to produce contextually
correct progressions on 10 held-out cart trajectories.


% ═══════════════════════════════════════════════════════════════════════════════
\section{Monitoring and Production Operations}
% ═══════════════════════════════════════════════════════════════════════════════

\subsection{Drift Detection}

Model performance is monitored on shadow traffic via:
\begin{itemize}
  \item NDCG on shadow prediction vs.\ training baseline (trigger retraining
        if NDCG drops $>$5\%)
  \item Segment-wise acceptance rate per tier/zone/age group (weekly report)
  \item Feature distribution shift (PSI $>$ 0.20 triggers feature re-engineering)
\end{itemize}

\subsection{Retraining Cadence}

The LightGBM model is retrained monthly on a 6-month rolling window of order
data. Promotion to production requires AUC $\geq 0.95$ on the latest holdout.
A blue-green deployment strategy ensures zero-downtime model swaps via a
\texttt{MODEL\_VERSION} environment variable.

\subsection{Known Limitations}

\begin{table}[h]
\caption{Known Limitations and Mitigations}
\centering
\begin{tabular}{p{2.2cm}p{2.4cm}p{2.4cm}}
\toprule
Limitation & Impact & Mitigation \\
\midrule
Synthetic data & Distribution gap vs.\ real transactions & Architecture is data-agnostic; replace with real data \\
\addlinespace
No restaurant features & Missing chain/quality signals & Add rating + chain flag as item features \\
\addlinespace
LLM explainer not deployed & Template-only explanations & Drop-in OpenAI/Gemini via \texttt{OPENAI\_API\_KEY} \\
\addlinespace
Monthly retraining & Drift between cycles & Shadow model + NDCG monitor catches drift early \\
\addlinespace
FP-Growth pool cap & Recall bounded by 50-item pool & Increase pool size; trade-off vs.\ latency \\
\bottomrule
\end{tabular}
\end{table}


% ═══════════════════════════════════════════════════════════════════════════════
\section{Related Work}
% ═══════════════════════════════════════════════════════════════════════════════

Add-on recommendation in e-commerce has been studied as complementary item
recommendation \cite{mcauley2015inferring, liu2017experimental}. The canonical
Amazon approach \cite{linden2003amazon} uses item-to-item collaborative
filtering, which succeeds at scale but ignores cart context.

Learning-to-Rank methods for recommendation are reviewed in
\cite{shi2010list}. LambdaRank \cite{burges2010ranknet} and its LambdaMART
variant (the precursor to LightGBM LambdaRank) are the most widely deployed
LTR models for production recommenders \cite{joachims2017}.

FP-Growth for recommendation pre-filtering is used in production at
Alibaba \cite{zhou2019deep} and Pinterest \cite{eksombatchai2018pixie},
though at larger scale with graph-based diffusion rather than item-count
co-occurrence.

Cold-start strategies for new users are surveyed in
\cite{lika2014facing}. Our three-tier fallback (warm/cool/cold) is a
pragmatic instantiation of standard practice \cite{schein2002methods}.

A/B testing methodology for recommendation systems is described in
\cite{kohavi2020trustworthy}. Our use of Welch's $t$-test with guardrail
metrics follows the framework advocated by \cite{deng2013improving}.


% ═══════════════════════════════════════════════════════════════════════════════
\section{Conclusion}
% ═══════════════════════════════════════════════════════════════════════════════

We presented CSAO, a production-ready, context-aware add-on
recommendation system achieving state-of-the-art ranking quality
(NDCG@8 = 0.876, Precision@8 = 0.382, AUC = 0.970) while operating at
P95 = 14.72\,ms---$20\times$ within the 300\,ms SLA. The system integrates
a segment-aware FP-Growth candidate generator, a LightGBM LambdaRank ranker
with 77 engineered features, a three-tier cold-start fallback, and a
production-benchmarked FastAPI server with full monitoring.

Two independent metric translation paths project a conservative $+22.1\%$
AOV lift ($+37.6\%$ upper bound), corresponding to \rupee5.157\,billion in
estimated annual incremental revenue at 200,000 daily orders. Segment-level
analysis reveals Tier-1 metro users and Millennials as highest-performing
cohorts, while Boomers and winter seasonality represent active improvement
opportunities for threshold tuning and reranking.

Future work includes: (1) deploying an LLM-based explainer for personalised
recommendation narratives; (2) extending to multi-restaurant session contexts
with cross-restaurant complementarity modelling; (3) adopting a streaming
feature store (Flink/Redis Streams) for real-time RFM signal updates;
and (4) exploring sequential models (SASRec, BERT4Rec) for users with
dense order histories on a GPU-equipped serving cluster.

% ── References ────────────────────────────────────────────────────────────────
\begin{thebibliography}{99}

\bibitem{ke2017lightgbm}
G.\ Ke, Q.\ Meng, T.\ Finley, T.\ Wang, W.\ Chen, W.\ Ma, Q.\ Ye, and
T.-Y.\ Liu, ``LightGBM: A highly efficient gradient boosting decision tree,''
\textit{Advances in Neural Information Processing Systems}, vol.\ 30, 2017.

\bibitem{burges2010ranknet}
C.\ J.\ C.\ Burges, ``From RankNet to LambdaRank to LambdaMART: An overview,''
\textit{Microsoft Research Technical Report}, MSR-TR-2010-82, 2010.

\bibitem{han2000mining}
J.\ Han, J.\ Pei, and Y.\ Yin, ``Mining frequent patterns without candidate
generation,'' in \textit{Proc.\ ACM SIGMOD International Conference on
Management of Data}, Dallas, TX, 2000, pp.\ 1--12.

\bibitem{joachims2017}
T.\ Joachims, A.\ Swaminathan, and T.\ Schnabel, ``Unbiased learning-to-rank
with biased feedback,'' in \textit{Proc.\ ACM WSDM}, Cambridge, UK, 2017.

\bibitem{craswell2008}
N.\ Craswell, O.\ Zoeter, M.\ Taylor, and B.\ Ramsey, ``An experimental
comparison of click position-bias models,'' in \textit{Proc.\ ACM WSDM},
Palo Alto, CA, 2008, pp.\ 87--94.

\bibitem{kang2018self}
W.-C.\ Kang and J.\ McAuley, ``Self-attentive sequential recommendation,''
in \textit{Proc.\ IEEE ICDM}, Singapore, 2018.

\bibitem{sun2019bert4rec}
F.\ Sun, J.\ Liu, J.\ Wu, C.\ Pei, X.\ Lin, W.\ Ou, and P.\ Jiang,
``BERT4Rec: Sequential recommendation with bidirectional encoder
representations from transformer,'' in \textit{Proc.\ CIKM}, Beijing, 2019.

\bibitem{grinsztajn2022tree}
L.\ Grinsztajn, E.\ Oyallon, and G.\ Varoquaux, ``Why tree-based models still
outperform deep learning on tabular data,'' in \textit{Advances in Neural
Information Processing Systems}, vol.\ 35, 2022.

\bibitem{linden2003amazon}
G.\ Linden, B.\ Smith, and J.\ York, ``Amazon.com recommendations:
Item-to-item collaborative filtering,'' \textit{IEEE Internet Computing},
vol.\ 7, no.\ 1, pp.\ 76--80, Jan./Feb.\ 2003.

\bibitem{mcauley2015inferring}
J.\ McAuley, R.\ Pandey, and J.\ Leskovec, ``Inferring networks of substitutable
and complementary products,'' in \textit{Proc.\ ACM KDD}, Sydney, 2015.

\bibitem{liu2017experimental}
Q.\ Liu, S.\ Wu, and L.\ Wang, ``An experimental comparison of complementary
and substitute based item recommendation,'' \textit{Neurocomputing}, vol.\ 250,
pp.\ 58--66, 2017.

\bibitem{mcinerney2018explore}
J.\ McInerney, B.\ Lacker, S.\ Hansen, K.\ Higley, H.\ Bouchard,
A.\ Gruson, and R.\ Mehrotra, ``Explore, exploit, and explain: Personalizing
explainable recommendations with bandits,'' in \textit{Proc.\ ACM RecSys},
Vancouver, 2018.

\bibitem{shi2010list}
Y.\ Shi, M.\ Larson, and A.\ Hanjalic, ``List-wise learning to rank with
matrix factorization for collaborative filtering,'' in \textit{Proc.\ ACM
RecSys}, Barcelona, 2010.

\bibitem{jannach2019measuring}
D.\ Jannach and M.\ Jugovac, ``Measuring the business value of recommender
systems,'' \textit{ACM Transactions on Management Information Systems},
vol.\ 10, no.\ 4, 2019.

\bibitem{kohavi2020trustworthy}
R.\ Kohavi, D.\ Tang, and Y.\ Xu, \textit{Trustworthy Online Controlled
Experiments: A Practical Guide to A/B Testing}. Cambridge University Press, 2020.

\bibitem{deng2013improving}
A.\ Deng, Y.\ Xu, R.\ Kohavi, and T.\ Walker, ``Improving the sensitivity of
online controlled experiments by utilising pre-experiment data,'' in
\textit{Proc.\ ACM WSDM}, Rome, 2013.

\bibitem{lika2014facing}
B.\ Lika, K.\ Kolomvatsos, and S.\ Hadjiefthymiades, ``Facing the cold start
problem in recommender systems,'' \textit{Expert Systems with Applications},
vol.\ 41, no.\ 4, pp.\ 2065--2073, 2014.

\bibitem{schein2002methods}
A.\ I.\ Schein, A.\ Popescul, L.\ H.\ Ungar, and D.\ M.\ Pennock,
``Methods and metrics for cold-start recommendations,'' in \textit{Proc.\ ACM
SIGIR}, Tampere, Finland, 2002.

\bibitem{ailawadi2014}
K.\ L.\ Ailawadi and P.\ W.\ Farris, ``Managing multi- and omni-channel
distribution: Metrics and research directions,'' \textit{Journal of
Retailing}, vol.\ 93, no.\ 1, pp.\ 120--135, 2017.

\bibitem{chintagunta2012}
P.\ K.\ Chintagunta, J.\ Chu, and J.\ Cebollada, ``Quantifying transaction
costs in online/offline grocery channel choice,'' \textit{Marketing Science},
vol.\ 31, no.\ 1, pp.\ 96--114, 2012.

\bibitem{zhou2019deep}
G.\ Zhou, X.\ Zhu, C.\ Song, Y.\ Fan, H.\ Zhu, X.\ Ma, Y.\ Yan, J.\ Jin,
H.\ Li, and K.\ Gai, ``Deep interest network for click-through rate
prediction,'' in \textit{Proc.\ ACM KDD}, London, 2018.

\bibitem{eksombatchai2018pixie}
C.\ Eksombatchai, P.\ Jindal, J.\ Z.\ Liu, Y.\ Liu, R.\ Sharma,
C.\ Sugnet, M.\ Ulrich, and J.\ Leskovec, ``Pixie: A system for
recommending 3+ billion items to 200+ million users in real-time,''
in \textit{Proc.\ WWW}, Lyon, France, 2018.

\bibitem{euromonitor2024}
Euromonitor International, ``Food delivery app retention benchmarks in
South/Southeast Asia,'' \textit{Passport: Consumer Foodservice}, 2024.

\end{thebibliography}

\end{document}
